\section{同一个十六进制数为什么能指向不同入口}

项目日志中会同时出现 \code{0x20}、\code{0x23} 和 \code{0x80}。若把所有
十六进制数都称为“地址”，那么 IRQ0 到来、加载用户代码段和执行系统调用就会像
三次含义不明的跳转。实际上，这些数字分别由设备端口译码、描述符表和中断描述符表
解释。数值本身没有携带类型，产生它的指令和读取它的硬件状态才决定命名空间。

\begin{keyidea}
看到一个关键数字时，必须同时写出\emph{产生者、解释者和命名空间}。
\code{0x20} 作为物理或虚拟地址、I/O 端口号、中断向量、设备命令字时可以同时
存在，彼此不会自动冲突。只有把一个命名空间的数字交给另一个解释者，才会形成故障。
\end{keyidea}

\subsection{先分开五种不会自动互换的数字}

\begin{longtable}{|L{0.18\linewidth}|L{0.22\linewidth}|L{0.25\linewidth}|L{0.25\linewidth}|}
\hline
数字的角色 & 典型产生位置 & 谁解释 & 是否经过页表 \\ \hline
虚拟地址 &
\code{load/store}、取指或指针运算 &
当前 CR3 下的 MMU 与页表 &
是；翻译后才形成物理访问 \\\tablerowrule
I/O 端口号 &
\code{IN/OUT} 指令的立即数或 DX &
处理器的独立 I/O 空间与平台设备译码 &
否；它不是内存地址 \\\tablerowrule
中断向量 &
异常硬件、PIC/APIC 或 \code{INT n} 指令 &
处理器以它索引当前 IDT &
向量本身不翻译；读取 IDT 门时才访问内存 \\\tablerowrule
段选择子 &
CS、SS、任务寄存器或门描述符 &
处理器拆出表索引、TI 与 RPL，再读取 GDT/LDT &
选择子本身不是表的字节地址 \\\tablerowrule
命令字或状态位 &
软件写设备寄存器，或设备返回寄存器值 &
具体控制器内部命令与状态译码 &
否；含义只在对应寄存器内成立 \\ \hline
\end{longtable}

下图用同一个 \code{0x20} 展示这种区别。每行都是独立关系，不表示上排输出会流入
下一排。

\begin{pdfdiagram}[scale=0.94]
  \node[os state, text width=36mm] (mem_in) at (-50mm,36mm)
    {\code{load [0x20]}\\虚拟地址操作数};
  \node[os compute, text width=31mm] (mmu) at (0,36mm)
    {MMU 与页表\\翻译地址};
  \node[os memory, text width=36mm] (mem_out) at (50mm,36mm)
    {当前地址空间中\\含 \code{0x20} 的位置};

  \node[os state, text width=36mm] (port_in) at (-50mm,18mm)
    {\code{OUT 0x20, AL}\\I/O 端口号};
  \node[os compute, text width=31mm] (port_decode) at (0,18mm)
    {平台 I/O 译码\\不查页表};
  \node[os memory, text width=36mm] (port_out) at (50mm,18mm)
    {主 PIC\\命令寄存器};

  \node[os state, text width=36mm] (vector_in) at (-50mm,0)
    {PIC 给出向量\\\code{0x20}=32};
  \node[os compute, text width=31mm] (idt) at (0,0)
    {索引 IDT[32]\\读取 16 字节门};
  \node[os control, text width=36mm] (vector_out) at (50mm,0)
    {IRQ0 内核入口\\CS 与 RIP 来源};

  \node[os state, text width=36mm] (selector_in) at (-50mm,-18mm)
    {加载选择子\\\code{0x23}};
  \node[os compute, text width=31mm] (selector_decode) at (0,-18mm)
    {拆成 index=4\\TI=0，RPL=3};
  \node[os memory, text width=36mm] (selector_out) at (50mm,-18mm)
    {GDT[4]\\Ring 3 代码段};

  \node[os state, text width=36mm] (command_in) at (-50mm,-36mm)
    {写入字节值\\\code{0x20}};
  \node[os compute, text width=31mm] (command_decode) at (0,-36mm)
    {PIC 命令译码\\非特定 EOI};
  \node[os control, text width=36mm] (command_out) at (50mm,-36mm)
    {清除当前最高优先级\\in-service 状态};

  \draw[os bus] (mem_in.east) -- (mmu.west);
  \draw[os bus] (mmu.east) -- (mem_out.west);
  \draw[os bus] (port_in.east) -- (port_decode.west);
  \draw[os bus] (port_decode.east) -- (port_out.west);
  \draw[os bus] (vector_in.east) -- (idt.west);
  \draw[os bus] (idt.east) -- (vector_out.west);
  \draw[os bus] (selector_in.east) -- (selector_decode.west);
  \draw[os bus] (selector_decode.east) -- (selector_out.west);
  \draw[os bus] (command_in.east) -- (command_decode.west);
  \draw[os bus] (command_decode.east) -- (command_out.west);
\end{pdfdiagram}
\diagramnote{指令种类和硬件入口先确定命名空间，随后数字才有含义。同为
\code{0x20} 的端口号、向量和 EOI 命令可以在一次 IRQ0 处理中连续出现，但它们
没有共享同一坐标原点，也没有发生三次数据搬运。}

\subsection{为什么端口号不能当作物理地址}

x86 为 \code{IN/OUT} 保留独立的 16 位 I/O 端口空间。执行
\code{OUT 0x20, AL} 时，处理器发起的是端口事务；当前 CR3、VMA 和 PTE 都不参与
决定目标。执行 \code{store [0x20], AL} 时，处理器发起的却是内存事务，首先按
当前页表翻译虚拟地址。两条指令即使携带相同数字，也会沿不同硬件通路到达不同对象。

项目使用的端口来自 PC 兼容平台契约，而不是处理器根据设备名称自动计算：

\begin{longtable}{|L{0.21\linewidth}|L{0.25\linewidth}|L{0.24\linewidth}|L{0.20\linewidth}|}
\hline
端口范围 & 项目中的对象 & 为什么必须使用该范围 & 若只改软件常量 \\ \hline
\code{0x20/0x21}、\code{0xA0/0xA1} &
主、从 8259A PIC 的命令与数据端口 &
PC 兼容芯片组在这些端口译码传统 PIC &
访问不到 PIC，屏蔽、重映射和 EOI 均失效 \\\tablerowrule
\code{0x40}、\code{0x43} &
PIT 通道 0 数据与模式命令 &
兼容 8254 的固定端口接口 &
计数初值写入其他设备或空端口 \\\tablerowrule
\code{0x60}、\code{0x64} &
i8042 数据端口与状态/命令端口 &
键盘控制器按传统 PC 端口布局工作 &
状态轮询与扫描码读取不再对应同一控制器 \\\tablerowrule
\code{0x1F0..0x1F7}、\code{0x3F6} &
主 ATA 数据、LBA、命令与控制寄存器 &
主 ATA 兼容通道使用这两段端口 &
LBA 字节和命令会写到错误端口 \\\tablerowrule
\code{0x3F8..0x3FF} &
COM1 的 16550 兼容寄存器窗口 &
平台把第一串口译码在传统 COM1 基址 &
最早启动证据消失，等待状态也不再可信 \\ \hline
\end{longtable}

这些数值看起来零散，是因为它们保留了不同年代 PC 外设的兼容接口。兼容性的收益是
固件、操作系统和虚拟机可以在枚举服务尚未建立时就访问设备；代价是地址不连续、寄存器
宽度和访问协议各不相同。另一种设计可以把设备放进 MMIO 区，或由 PCI BAR 在启动时分配
地址，但软件必须先枚举设备、建立页表映射并选择正确缓存属性。当前项目选择传统端口，
是为了先闭合启动、异常和设备事件，不表示现代驱动都应使用这些固定值。

\subsection{为什么 \code{0x08}、\code{0x1B} 和 \code{0x23} 不是表地址}

段选择子的低三位不是地址低位。处理器按
\[
  \text{selector}=(\text{descriptor index}\ll3)
  \mathbin{|}(\text{TI}\ll2)\mathbin{|}\text{RPL}
\]
拆解它。项目当前只使用 GDT，所以 TI 为 0；RPL 保存请求特权级。由此可复算：

\begin{longtable}{|L{0.16\linewidth}|L{0.18\linewidth}|L{0.26\linewidth}|L{0.30\linewidth}|}
\hline
选择子 & 位拆解 & 指向的描述符 & 为什么是这个数 \\ \hline
\code{0x08} &
\(1\ll3\) &
GDT[1]，Ring 0 代码段 &
索引 1，RPL 0 \\\tablerowrule
\code{0x10} &
\(2\ll3\) &
GDT[2]，Ring 0 数据段 &
索引 2，RPL 0 \\\tablerowrule
\code{0x1B} &
\((3\ll3)|3\) &
GDT[3]，Ring 3 数据段 &
索引 3，加 RPL 3 后低位成为 \code{11b} \\\tablerowrule
\code{0x23} &
\((4\ll3)|3\) &
GDT[4]，Ring 3 代码段 &
索引 4，加 RPL 3；不是物理地址 \\\tablerowrule
\code{0x28} &
\(5\ll3\) &
GDT[5..6]，64 位 TSS 描述符 &
TSS 描述符占两个表项，选择子指向第一个 \\ \hline
\end{longtable}

GDT[0] 保持全零，使空选择子不可能偶然命中一个可用段。表项顺序是项目策略，
不是 x86 要求的唯一顺序；换一种顺序完全可行，但 GDT 构造、门描述符、汇编常量、
用户返回帧和运行时验证必须一起更新。只改 \code{0x23} 会让处理器读取另一个表项，
随后可能因类型或特权不匹配产生 \#GP。

在 64 位模式下，普通代码段和数据段不再承担传统平坦寻址的全部基址与界限工作，
但 CS 中的特权级、代码段类型和长模式属性仍参与入口与返回检查。选择子因此不是
“历史遗留但可以随便填的数字”；它仍把一次控制转移绑定到具体权限契约。

\subsection{为什么异常占 \code{0..31}，PIC 从 \code{32} 开始}

向量是 IDT 的表项索引。x86-64 的一个 IDT 门占 16 字节，因此向量 \(v\) 对应的
门位于
\[
  \text{IDTR.base}+16v.
\]
这里发生的是“用索引读取门描述符”，不是把向量当作入口 RIP。门描述符再给出
目标偏移、代码段选择子、门类型、DPL 和可选 IST。

向量 \code{0..31} 由处理器架构保留给除零、无效操作码、双重故障、页故障等异常。
传统 PIC 上电后的默认映射会与这个范围重叠，因此内核把主 PIC 的八条 IRQ 映射到
\code{32..39}，从 PIC 的八条 IRQ 映射到 \code{40..47}。基址按 8 对齐，是因为
8259A 在 8086 模式下由初始化字给出向量高位，低三位由八条 IRQ 编号补入。

若仍让 IRQ0 使用向量 8，处理器进入 IDT[8] 时，统一现场无法仅凭向量判断这是 PIT
事件还是双重故障。把 PIC 起点改为 32 后，架构异常与设备事件拥有不相交的入口范围：
\[
  \text{IRQ vector}=32+\text{IRQ number}.
\]
映射到别的、同样空闲且满足对齐的范围也可以，但 PIC 初始化、IDT 门、分发器、测试
和调试符号必须采用同一基址。

\subsection{为什么系统调用先使用 \code{INT 0x80}}

\code{0x80}=128 位于架构异常和传统 PIC 范围之外，项目把 IDT[128] 建成 DPL3
interrupt gate。DPL3 是允许 Ring 3 通过 \code{INT 0x80} 显式进入的关键；
若仍为 DPL0，用户指令会先得到 \#GP，根本到不了系统调用分发器。其余只允许内核
触发的软件门不能因为方便而全部改成 DPL3，否则用户程序也能合成那些入口。

向量 128 并非 x86-64 强制的系统调用号码。项目沿用熟悉的 \code{INT 0x80}
约定，并复用已经验证的 IDT、统一异常帧和 TSS.RSP0，在首次用户边界中减少同时新增
的机制。换成另一个未占用向量可行，但用户汇编、共享 ABI、IDT 和入口校验必须同步；
直接使用 \code{SYSCALL/SYSRET} 也可行，却还要配置 EFER.SCE、STAR、LSTAR 和
FMASK，并在入口自行取得可信内核栈。后者延迟更低，但会把新的快速入口契约加入同一
演进阶段。

\subsection{沿一次 IRQ0 把所有数字放回各自解释者}

设 PIT 到期并经主 PIC 的 IRQ0 请求处理器。完整路径可以逐步复算：

\begin{enumerate}
  \item PIT 只拉起 IRQ0；此时还没有 IDT 向量，也没有内核入口地址。
  \item PIC 用已配置的基址计算 \(32+0=\code{0x20}\)，向处理器提供中断向量。
  \item 处理器读取 \(\text{IDTR.base}+32\times16\) 处的 IDT 门。读取表项不会
        搬动处理器即将执行的代码，只取得新的 CS、RIP 和门属性。
  \item 门中的 \code{0x08} 按选择子规则指向 GDT[1]；处理器据此进入 Ring 0
        代码段，再从门内偏移开始取指。
  \item 内核完成 PIT 记账后，把命令字 \code{0x20} 写入 I/O 端口
        \code{0x20}，通知主 PIC 结束当前 in-service 请求。
\end{enumerate}

最后一步同时出现两个 \code{0x20}：前一个是 EOI 命令字，后一个是 PIC 命令端口。
更早的 \code{0x20} 则是 IDT 向量。三者数值相等只是协议碰巧重合；它们分别由 PIC
命令译码、平台端口译码和处理器 IDT 逻辑解释。这里没有“同一地址变化三次”，也没有
数据被反复复制。

\begin{warningbox}
调试入口故障时，先给每个数字补上类型。把端口号交给内存读写，会访问当前页表中的
另一位置；把选择子直接乘表项大小而忽略 RPL，会定位错误表项；把 IRQ 号直接当向量，
会进入架构异常范围；把向量当 RIP，则会跳到一个完全没有由门验证的地址。只有解释者
和命名空间同时正确，十六进制数才构成完整事实。
\end{warningbox}
